\paragraph{Data.} Classification uses five public CT cohorts with \numLabelsN{} binary labels (Table~\ref{tab:cohorts} and Figure~\ref{fig:examples} in Appendix~\ref{app:cohorts}): the CVPR26 cohort, the release of the CVPR 2026 3-D CT foundation-model challenge, which groups AMOS \citep{ji2022amos} with 15 disease labels, COVID-CTset \citep{rahimzadeh2021covidctset} and LUNA25 nodule malignancy \citep{peeters2026luna25}; FDG-PET/CT \citep{gatidis2022fdgpetct}, whose four diagnosis groups we predict from the CT alone; RSNA 2022 cervical-spine fractures \citep{lin2023rsnacsfdata}; RSNA 2023 abdominal trauma \citep{rudie2024ratic}; and RSNA 2019 intracranial hemorrhage \citep{flanders2020rsnaich}. Segmentation uses AMOS22 (15 organs; \citealp{ji2022amos}), KiTS23 \citep{heller2023kits21}, LiTS \citep{bilic2023lits}, MSD Pancreas \citep{antonelli2022msd} and TotalSegmentator \citep{wasserthal2023totalsegmentator,wasserthal2026totalsegdata} with its 117 structures merged into 15 (Table~\ref{tab:totalseg}); \numStructuresN{} structures in all. For the CVPR26 cohort and AMOS22 the official validation cases are divided into a selection and a held-out half; the other sources use fixed seeded splits of the released labeled cases, and TotalSegmentator holds out its official test split. VoCo is run on FDG-PET/CT, RSNA 2022, RSNA 2023 and all segmentation datasets. Several evaluation sets appear, as images, in the pretraining data of SPECTRE and VoCo (Table~\ref{tab:overlap}).

\paragraph{Optimization.} All runs use AdamW with 10\% linear warm-up, cosine decay to $10^{-6}$ and bf16 precision. A classification run trains 8000 steps at batch size 1 on a masked, positive-weighted binary cross-entropy; a segmentation run trains 8000 steps of four crops on soft Dice plus binary cross-entropy; two-stage methods split the budget in halves. A joint model thus uses $1/K$ of the steps of the $K$ per-task models it replaces, except the 15-structure joint segmentation models (AMOS22, TotalSegmentator), which train 24\,000 steps. Each method has one default learning rate, set by what it trains: $10^{-3}$ for the classification linear probe and the segmentation decoder, $10^{-2}$ for the segmentation linear probe, $10^{-4}$ for the attentive probe and the adapters, and $2\times10^{-5}$ for encoder weights; none was tuned on our data (Tables~\ref{tab:optim} and~\ref{tab:optimseg}). The checkpoint with the best selection macro AUROC or mean Dice, evaluated every 1000 steps, is kept and its held-out metric reported.

\paragraph{Evaluation.} AUROC per label and Dice per structure and case are the metrics \citep{maierhein2024metrics}. Methods are compared per label--mode pair with DeLong's test \citep{delong1988} or a paired case bootstrap on Dice, with Benjamini--Hochberg correction \citep{benjamini1995fdr}, and ranked within each pair that has every method. The main comparisons are summarized at label level, each label counting once with its two modes averaged, with 95\% intervals from a bootstrap over cohorts, labels and held-out cases (Appendix~\ref{app:stats}); as four or five cohorts per backbone make the cohort stage optimistic, we also count the cohorts with a positive mean difference. VoCo, with three of the five classification cohorts, is reported per backbone and never pooled. A method whose held-out predictions barely vary between cases did not train: SPECTRE Full FT, with a median standard deviation of held-out probabilities of \numSpreadSpectreFullft{} against at least \numSpreadOthermin{} for every other method; it stays in the tables and out of ranks and averages over methods. The joint models of Lin, LoRA, Attn-only FT and Full FT (classification) and of Dec, LoRA and Full FT (segmentation) were trained with two further seeds (Appendix~\ref{app:revision}).
